%%%%%%%%%%%%%%%%%%%%%%%%%%%%%%%%%%%%%%%%%%%%%%%%%%%%%%%%%%%%%%%%%%%%%%%%%%%%%%%
% IE 692 -- Course Project Proposal
%
% Compile:   latexmk -pdf proposal.tex      (or: pdflatex proposal.tex, twice)
% Clean up:  latexmk -c
%
% Switch \usepackage[draft]{ie692proposal} to [final] to hide every \todo{}
% and \note{} before you submit.
%
% WHAT THE HANDOUT ACTUALLY REQUIRES OF THE PROPOSAL (Submission Instructions):
%   (1) a single typed PDF listing the team members,
%   (2) an outline of the three proposed enhancements,
%   (3) confirmation that all members can present Tue Nov 24, 12:55-2:10pm EST.
% Everything else below is optional, and is here because the handout says
% feedback is given after the proposal -- more detail buys better feedback.
%%%%%%%%%%%%%%%%%%%%%%%%%%%%%%%%%%%%%%%%%%%%%%%%%%%%%%%%%%%%%%%%%%%%%%%%%%%%%%%

\documentclass[11pt]{article}

\usepackage[draft]{ie692proposal}   % <- change to [final] before submitting
\input{math_commands.tex}

\coursename{IE 692: Optimization Methods for Reinforcement Learning}
\doctitle{Course Project Proposal}
\proposaldue{Fri, Sep 18th (5pm EST)}
\finaldue{Tue, Nov 24th (5pm EST)}

\begin{document}

\makeheaderbox

%%%%%%%%%%%%%%%%%%%%%%%%%%%%%%%%%%%%%%%%%%%%%%%%%%%%%%%%%%%%%%%%%%%%%%%%%%%%%%%
\section{Team and Logistics}
%%%%%%%%%%%%%%%%%%%%%%%%%%%%%%%%%%%%%%%%%%%%%%%%%%%%%%%%%%%%%%%%%%%%%%%%%%%%%%%
% REQUIRED: team members, and the presentation-availability confirmation.

\noindent
\begin{tabular}{@{}lll@{}}
\toprule
Name & Email & Attendance \\
\midrule
{Jack Morris} & {jmorr185@vols.utk.edu} & In-person \\
{Anjila Budathoki} & {abudatho@vols.utk.edu} & In-person \\
{Manish Dhakal} & \ {mdhakal1@vols.utk.edu} & In-person \\
\bottomrule
\end{tabular}

\vspace{0.5\baselineskip}
\noindent
\textbf{All team members confirm that they are available to present on Tuesday,
November 24th from 12:55pm to 2:10pm EST.}


% %%%%%%%%%%%%%%%%%%%%%%%%%%%%%%%%%%%%%%%%%%%%%%%%%%%%%%%%%%%%%%%%%%%%%%%%%%%%%%%
\section{Proposed Enhancements}
% %%%%%%%%%%%%%%%%%%%%%%%%%%%%%%%%%%%%%%%%%%%%%%%%%%%%%%%%%%%%%%%%%%%%%%%%%%%%%%%


% We will implement a tabular RL library from scratch and apply it to the two
% required environments: the simplified blackjack game and a $10\times10$
% GridWorld with 10 hazards, 3 goal states, and 1 origin state, both with
% discount factor $\disc = 0.9$. Our baseline consists of on-policy Monte Carlo
% (OMC) for policy evaluation and stochastic policy mirror descent (SPMD) for
% policy optimization. On top of this baseline 
We propose three enhancements:
(i) SPMD with three Bregman divergences, (ii) True Online TD(\texorpdfstring{$\lambda$}{lambda}) Evaluation, and (iii) entropy-regularized SPMD.



\begin{enhancement}{SPMD with Three Bregman Divergences}{Policy optimization}
\begin{algorithm}
\caption{Stochastic Policy Mirror Descent}
\label{alg:spmd}
\begin{algorithmic}[1]
\Procedure{Optimization}{$\niter, \polinit$}
\State $\polk \gets \polinit$
\For{$\iter \gets 0$ \To $\niter$}
    \State $\Qtilk \gets \Qtilk[\state,\action] \approx \Qpi[\state,\action] \allsa$
    \State $\polk \gets \underset{\polof[\polk]{\state} \in \simplexA}{\argmin}
      \set*{ \dotp{\Qtilk[\state,\cdot]}{\polof{\state}} \stepk
             + \bregD{\polof[\polk]{\state}}{\polof{\state}} }$
\EndFor
\Return $\polk$
\EndProcedure
\end{algorithmic}
\end{algorithm}


\field{Description}
We instantiate SPMD with three Bregman divergences generated by different
distance-generating functions: squared Euclidean distance, KL divergence, and Tsallis divergence.
For example, when KL divergence is used as $\bregD{\polof[\polk]{\state}}{\polof{\state}}$ (In Algorithm \ref{alg:spmd}) we get

\[
    \underset{p \in \simplexA}{\argmin}
      \set*{ \dotp{\Qtilk[\state,\cdot]}{p} \stepk
             + \sum_{\action \in \cA} p_a \ln \frac{p_a}{\polk (\state,\action)} }
\]             
as our update to $\polk$. 


\field{Motivation} We aim to measure the convergence and performance of the learned $\polk$ under each divergence for each MDP problem.  Where there is an observed difference in either metric, we will identify empirically or derive its source.  



\field{Reference}
Reinforcement Learning through Optimization, Section 2.1.1 (Bregman's divergence) and Section 2.2 (Policy
mirror descent). 

% \field{Hypothesis}
% \todo{state it}
% The divergence sets the geometry in which the update is taken, so we expect the
% KL and Tsallis updates --- which respect the simplex --- to outperform the
% Euclidean projection on GridWorld, where near-deterministic optimal actions
% push iterates toward the boundary of $\simplexA$. We expect smaller differences
% on blackjack, where the optimal policy is less concentrated.

% \field{Parameters to tune}
% Step size $\stepk$; number of rollouts per state--action pair $\rollouts$;
% Tsallis entropic index $q$.

% \field{Implementation sketch}
% A single \texttt{BregmanUpdate} interface with three subclasses, so the SPMD
% loop is written once. Euclidean uses a sort-based simplex projection; Tsallis
% uses bisection on the normalizer with a tolerance of $10^{-8}$.
\end{enhancement}

\begin{enhancement}{True Online TD(\texorpdfstring{$\lambda$}{lambda}) Evaluation}{Policy evaluation}
\field{Description}

The standard On-Policy Monte Carlo  evaluator estimates the action-value
function under each fixed SPMD policy using the complete sampled return
\[
G_t
=
\sum_{j=t}^{T-1} \gamma^{\,j-t} c_j,
\]
with the update
\[
\Qtil(s_t,a_t)
\leftarrow
\Qtil(s_t,a_t)
+
\stepsize\bigl(G_t-\Qtil(s_t,a_t)\bigr).
\]
and updates the estimate toward $G_t$. We use True Online TD($\lambda$),
which updates the action-value estimate $\Qtil(s,a)$ after every transition using
the temporal-difference error
\[
\delta_t
=
c_t
+\gamma
\sum_{a' \in \cA}
\polk(a' \mid s_{t+1}) \Qtil(s_{t+1},a')
-\Qtil(s_t,a_t).
\]
Monte Carlo, which relies on complete sampled returns and can therefore produce
delayed and relatively high-variance updates, while, True Online TD($\lambda$) learns
incrementally while using eligibility traces to propagate TD errors to recently
visited state-action pairs. This allows delayed goal and hazard costs to
influence earlier estimates without waiting for a complete episode.

\field{Motivation} The tuned On-Policy Monte Carlo baseline relies on complete sampled returns, which delays value updates until an episode is completed and can produce relatively high-variance estimates. The proposed enhancement learns incrementally from each transition while using eligibility traces to propagate TD errors to recently visited state-action pairs. This allows delayed goal and hazard costs to influence earlier estimates without waiting for a complete episode. 

% The true online correction maintains equivalence with the corresponding online forward view as the estimates change, while $\lambda$ controls the balance between short-term bootstrapping and longer-horizon returns.

% \field{Parameters to tune}
% Learning rate $\stepsize$, $\lambda$, and the number of transitions used in each policy-evaluation call.

% \field{Implementation sketch}
% Under each fixed SPMD policy, $\Qtil$ is updated after every transition using True Online TD, with eligibility traces propagating information to recently visited state-action pairs. We aim to measure performance by comparing with the On-Policy Monte Carlo baseline using convergence, variance, and sample efficiency.
\end{enhancement}

\field{Reference}

\begin{enumerate}[leftmargin=2em, topsep=0pt, itemsep=0pt, parsep=0pt]
    \item H. van Seijen and R. Sutton, ``True Online TD($\lambda$),'' ICML, 2014.
    \item Policy Evaluation, Section 3.2.3 (TD(0), TD($\lambda$), and Other Related Methods).
\end{enumerate}


\begin{enhancement}{Entropy-Regularized SPMD}{Policy optimization}
\field{Description}
We add a negative-entropy regularizer $\entcoef\,\negent(p)$, $\negent(p) =
\sum_{\action} p_{\action} \log p_{\action}$, to the SPMD subproblem (2.2.21),
\[
  \polknext(\cdot \mid \state)
  = \argmin_{p \in \simplexA}\;
    \stepk\big[\inner{\Qtil[\state,\cdot],\, p} + \entcoef\,\negent(p)\big]
    + \KL{p}{\polk(\cdot \mid \state)},
\]
which is the entropy-regularized problem of Section 2.1.2. With the KL
divergence the update stays closed-form,
\[
  \polknext(\action \mid \state)
  \;\propto\;
  \polk(\action \mid \state)^{\frac{1}{1+\stepk\entcoef}}
  \exp\!\left(-\frac{\stepk\,\Qtil[\state,\action]}{1+\stepk\entcoef}\right),
\]
and $\entcoef = 0$ recovers baseline SPMD (2.2.3). The regularizer pulls the
policy toward the Boltzmann policy $\pol \propto \exp(-\Qpi/\entcoef)$, so
every action keeps a probability bounded away from zero.

\field{Reference}
Course text, Sections 2.1.2, 2.2.1, and 2.3.1.

\field{Motivation}
Online policy evaluation needs every action to keep positive probability
(2.3.5), but unregularized SPMD converges toward a deterministic policy, so
rarely chosen actions stop being sampled and their value estimates go stale.
On GridWorld this can lock the policy into a route that never reaches a goal.
Entropy regularization slows this collapse, at the price of converging to the
regularized optimum. We therefore expect a non-monotone effect in $\entcoef$,
and will also test an annealed schedule $\entcoef_k = \entcoef_0/(k+1)^{p}$
that keeps early exploration without the bias.

% \field{Parameters to tune}
% $\entcoef$ (or $\entcoef_0$ and $p$ when annealed); step size $\stepk$;
% rollouts per iteration $\rollouts$.

% \field{Implementation sketch}
% A flag on the KL mirror step, computed in log space to avoid underflow. We
% report the unregularized $\Vpik[\state_0]$ so that curves across $\entcoef$
% are comparable, and additionally track $\min_{\state,\action} \polk(\action
% \mid \state)$ and GridWorld state coverage per iteration.
\end{enhancement}



%%%%%%%%%%%%%%%%%%%%%%%%%%%%%%%%%%%%%%%%%%%%%%%%%%%%%%%%%%%%%%%%%%%%%%%%%%%%%%%
% \section{Experimental Plan}
% %%%%%%%%%%%%%%%%%%%%%%%%%%%%%%%%%%%%%%%%%%%%%%%%%%%%%%%%%%%%%%%%%%%%%%%%%%%%%%%
% % OPTIONAL but cheap. Shows you have read the "project design and execution"
% % paragraph of the handout, which is where the 50 points live.

% Following the handout, every method will be measured by (i) the estimated
% state-value function $\Vpit[\state_0]$ at the origin state versus iteration $t$
% and (ii) the number of observations $\nobs$ consumed at each iteration. For
% each enhancement we will grid-search the three parameters listed above over
% three values each ($27$ combinations), scoring each combination on a reduced
% budget of \todo{5--10\%} of the full observation budget $\budget$, then run the
% best combination for the full budget.

% \begin{center}
% \begin{tabular}{@{}lll@{}}
% \toprule
% Enhancement & Parameters & Candidate values \\
% \midrule
% Bregman divergences & $\stepsize$, $\rollouts$, $q$ & \todo{fill in} \\
% $\lambda$-advantages & $\lambda$, $\rollouts$, $\stepsize$ & \todo{fill in} \\
% Entropy regularization & $\entcoef$, $\stepsize$, $\rollouts$ & \todo{fill in} \\
% \bottomrule
% \end{tabular}
% \end{center}




%%%%%%%%%%%%%%%%%%%%%%%%%%%%%%%%%%%%%%%%%%%%%%%%%%%%%%%%%%%%%%%%%%%%%%%%%%%%%%%
% \section{References, Collaboration, and AI Use}
% %%%%%%%%%%%%%%%%%%%%%%%%%%%%%%%%%%%%%%%%%%%%%%%%%%%%%%%%%%%%%%%%%%%%%%%%%%%%%%%
% % REQUIRED BY COURSE POLICY. The handout: AI use must be cited AND must show
% % the prompts and generated output; collaborators must be named with a
% % high-level description of what was discussed. An AI link or screenshot
% % without prompts and output is treated as plagiarism.

% \noindent\textbf{References.} \todo{course text sections, any papers}

% \vspace{0.3\baselineskip}
% \noindent\textbf{Collaboration.} \todo{names of students outside the team and a
% high-level description of what was discussed, or ``none''}

% \vspace{0.3\baselineskip}
% \noindent\textbf{AI use.} \todo{describe which tool was used and for what;
% attach or link a transcript showing the prompts and the generated output, per
% the course policy}

\end{document}
